% LỜI TỰA --- Quyển VIII
\chapter*{Lời Tựa}
\addcontentsline{toc}{chapter}{Lời Tựa}
\markboth{Lời Tựa}{Lời Tựa}

\begin{truyen}{Lời Tựa}{Rough but Real: Student Stories}
\chuhoa{G}{iáo viên bắt đầu bộ sách này với một câu hỏi đơn giản: làm thế nào để mỗi buổi học có thể để lại thứ gì đó đáng nhớ hơn là bài thi?}
\textit{(A teacher began this series with a simple question: how can each class leave behind something more memorable than a test?)}

Những câu chuyện hay nhất về học sinh không phải là những câu chuyện ngọt ngào.
\textit{(The best stories about students are not the sweet ones.)}

Quyển tám chọn những chuyện thô ráp, trực tiếp, đôi khi khó nghe. Về điểm số, về áp lực, về so sánh, về bắt nạt, về mạng xã hội, về việc chọn nghề và về nỗi sợ thất bại. Những chuyện đó không được làm đẹp lên~--- chúng được kể thật, vì thật mới chạm được vào người đọc.
\textit{(Volume eight chooses the rough, direct, sometimes hard-to-hear stories. About grades, pressure, comparison, bullying, social media, choosing a career, and the fear of failure. These stories are not made pretty~--- they are told truthfully, because truth is what reaches the reader.)}

Bộ sách <<Truyện Tích Cảm Hứng Song Ngữ>> gồm 24 quyển, mỗi quyển một chủ đề riêng. Quyển VIII này mang chủ đề: \textbf{điểm số, áp lực, so sánh, bắt nạt và học đường thật}.
\textit{(The series ``Inspiring Bilingual Tales'' contains 24 volumes, each with its own theme. Volume VIII focuses on: \textbf{điểm số, áp lực, so sánh, bắt nạt và học đường thật}.})
\end{truyen}

\begin{baihoc}
Câu chuyện thật không cần phải đẹp~--- nó chỉ cần đủ dũng cảm để được kể.
\textit{(A true story does not need to be beautiful~--- it only needs the courage to be told.)}
\end{baihoc}

\begin{ghinhoanh}
\textbf{Short English to remember:}

Honesty in storytelling creates real connection.

Real student life is rough, and that is okay.
\end{ghinhoanh}

\begin{tuhocnhanh}
1. Câu chuyện nào trong quyển này khiến em cảm thấy mình được hiểu? \textit{(Which story in this volume made you feel understood?)}

2. Áp lực học đường nào em đang chịu mà chưa nói với ai? \textit{(What school pressure are you carrying that you haven't shared with anyone?)}

3. Nếu viết một câu chuyện học đường thật của mình, em sẽ kể chuyện gì? \textit{(If you wrote one true story from your own school life, what would it be?)}
\end{tuhocnhanh}

\begin{center}
\textit{\color{truyengold}``Thô nhưng thật còn hơn đẹp mà rỗng.''}
\end{center}

\begin{flushright}
\textbf{Tôi Nguyễn Văn Sang}\\
\textit{GV THPT Nguyễn Hữu Cảnh - TP HCM}
\end{flushright}

\ngancach
